%Title: Beamer Presentation Template
%Author: LISA
%Year: 2020

\documentclass[10pt,aspectratio=169]{beamer}

%
%Setting file
%

\usepackage[T1]{fontenc}
\usepackage[utf8]{inputenc}

\usepackage[english]{babel}

\usepackage{graphicx}
\graphicspath{{images/}}
\usepackage{float}
\usepackage{tikz}
\usepackage{caption}
\usepackage{subcaption}
\usepackage[absolute,overlay]{textpos}
\usepackage{xcolor}

\usetheme{default}
\usefonttheme{structurebold}


% ---------------------------------
% color definitions
\usepackage{color}
\definecolor{LISA_BLUE}{cmyk}{0.99,0.88,0.29,0.18}
\setbeamercolor{normal text}{fg=LISA_BLUE}
\setbeamercolor{frametitle}{fg=LISA_BLUE}

\colorlet{itemizecolor}{red}

\newcommand\insertlocation{}  % Empty by default.
\newcommand\location[1]{\renewcommand\insertlocation{#1}}

\newcommand\insertcoauthor{}  % Empty by default.
\newcommand\coauthor[1]{\renewcommand\insertcoauthor{#1}}

\newcommand\insertperiod{}  % Empty by default.
\newcommand\period[1]{\renewcommand\insertperiod{#1}}



\setbeamertemplate{itemize items}[circle]
\setbeamercolor{title}{fg=white}

\setbeamertemplate{title page}
{
	\begin{textblock*}{0.7\paperwidth}(0.03\paperwidth,0.25\paperheight)
         

      \usebeamerfont{title}\textcolor{white}{\inserttitle}%
      \ifx\insertsubtitle\@empty%
      \else%
        \vskip0.25em%
        {\usebeamerfont{subtitle}\usebeamercolor[fg]{subtitle}\insertsubtitle\par}%
      \fi%     

      \vspace{0.1\textheight}      
      \usebeamerfont{author}\textcolor{white}{\insertauthor}
      \vskip1em%
      \usebeamerfont{institute}\textcolor{white}{\insertinstitute}

  \end{textblock*}
  
}

\makeatletter
    \newenvironment{SectionTitle}{
        % \setbeamertemplate{headline}[default]
        \setbeamercolor{normal text}{fg=LISA_BLUE}
        \setbeamercolor{frametitle}{fg=LISA_BLUE}
        \setbeamercolor{title}{fg=LISA_BLUE}
        \setbeamercolor{normal text}{fg=LISA_BLUE}
        \usebeamercolor[fg]{normal text}
        \usebeamercolor[fg]{footline}
        \usebackgroundtemplate%
        {%
            \includegraphics[width=\paperwidth,height=\paperheight]{Section.pdf}%
        }

    }{}
\makeatother

\useoutertheme{smoothbars} 
\setbeamertemplate{footline}[frame number]{}
\setbeamertemplate{navigation symbols}{}
\setbeamertemplate{footline}{}


\usebackgroundtemplate%
{%
    \includegraphics[width=\paperwidth,height=\paperheight]{Slide.pdf}%
}


\setbeamertemplate{footline}
{
 \begin{beamercolorbox}[wd=\paperwidth, ht=0.05\paperheight, dp=5.2pt,leftskip=2cm,rightskip=.3cm]{section in head/foot}
 \hspace{0.095cm}\textcolor{white}{\insertdate \hfill  \insertframenumber}\phantom{x}\vskip2pt
 \end{beamercolorbox}%
 \ifnum \theframenumber=0
        % This is the title frame, do nothing

    \else
        % These would be all the following frames
        % <<<Your template code goes here>>>
        \begin{textblock*}{0.3\paperwidth}(0.13\paperwidth,0.975\paperheight)
           \textcolor{white}{\insertlocation{}}
       \end{textblock*}
    \fi
}

\makeatletter
\setbeamertemplate{frametitle}{
  \vspace{0.05\textheight}
  \begin{beamercolorbox}[center]{frametitle}
      \usebeamerfont{frametitle} \insertframetitle%
  \end{beamercolorbox}%
}
\makeatother
% \beamertemplatenavigationsymbolsempty
% \setbeamercovered{transparent}

\setbeamertemplate{headline}
{%

}

%-----------------------------------------Title page settings-----------------------------------------%
\title{The Physics Behind Colors}
\subtitle{}
\author{Andrea Raggio}
\institute{19:00 09/12/21\\Tuopillinen Tiedett\"{a} @ Vihre\"{a} Haltiatar}
% \period{April-2020}
% \location{Pint of Science}


%-----------------------------------------Title page settings-----------------------------------------%
\usepackage{animate}
\usepackage{mathtools}


\begin{document}

{
  \usebackgroundtemplate{\includegraphics[width=\paperwidth]{Title.pdf}}
	\begin{frame}[noframenumbering, plain]
		\titlepage
	\end{frame}
}

\begin{SectionTitle}
{
  \usebackgroundtemplate{\includegraphics[width=\paperwidth]{Section_light_off.pdf}}
	\begin{frame}[noframenumbering, plain]
		\centering
		\vspace{0.1\textheight}
		% \textbf{\LARGE LIGHT}
	\end{frame}
}
{
  \usebackgroundtemplate{\includegraphics[width=\paperwidth]{Section_light_on.pdf}}
	\begin{frame}[noframenumbering, plain]
		\centering
		\vspace{0.1\textheight}
		% \textbf{\LARGE LIGHT}
		\vspace{0.1\textheight}
		% \textbf{\LARGE LIGHT}
	\end{frame}
}
\end{SectionTitle}
\begin{frame}{What is light}
	\begin{textblock*}{0.9\paperwidth}(0.05\paperwidth,0.2\paperheight)
		\centering
		Let's start from a broader perspective $\rightarrow$ \textbf{electromagnetic radiation (EMR)}
	\end{textblock*}
	\begin{center}
		\only<2>{\animategraphics[loop,controls,width=0.5\linewidth]{20}{E_Wave/E_wave-}{0}{35}}
		\only<3>{\animategraphics[loop,controls,width=0.5\linewidth]{20}{EM_Wave/EM_wave-}{0}{35}}
	\end{center}

	% \only<1>{\includegraphics[width=\textwidth]{EM_wave_1}}%
	% \only<2>{\includegraphics[width=\textwidth]{EM_wave_2}}%
	% \only<3>{\includegraphics[width=\textwidth]{EM_wave_3}}%
\end{frame}

\begin{frame}{EM spectrum}
	\centering
	Light is a small portion of the EM spectrum
	\only<1>{\includegraphics[width=\textwidth]{EM_spectrum-1.pdf}}%
	\only<2>{\includegraphics[width=\textwidth]{EM_spectrum-2.pdf}}%
\end{frame}

\begin{frame}{Sun spectrum}
	\centering
	\only<1>{\includegraphics[width=\textwidth]{sun.png}}%
	\only<2>{\includegraphics[width=0.7\textwidth]{sun.jpg}}
\end{frame}
\begin{frame}{Elements spectra}
	\centering
	\includegraphics[width=0.45\textwidth]{A-spectra-H-He-Li-Be.png}
	\includegraphics[width=0.45\textwidth]{E-spectra-H-He-Li-Be.png}
	\begin{textblock*}{0.05\paperwidth}(0.02\paperwidth,0.3\paperheight)
		\centering
		\textbf{H}
	\end{textblock*}	
	\begin{textblock*}{0.05\paperwidth}(0.02\paperwidth,0.43\paperheight)
		\centering
		\textbf{He}
	\end{textblock*}	
	\begin{textblock*}{0.05\paperwidth}(0.02\paperwidth,0.56\paperheight)
		\centering
		\textbf{Li}
	\end{textblock*}	
	\begin{textblock*}{0.05\paperwidth}(0.02\paperwidth,0.690\paperheight)
		\centering
		\textbf{Be}
	\end{textblock*}
	
\end{frame}
\begin{SectionTitle}
	{
	  \usebackgroundtemplate{\includegraphics[width=\paperwidth]{Section.pdf}}
		\begin{frame}
			\begin{textblock*}{0.5\paperwidth}(0.25\paperwidth,0.45\paperheight)
				\centering
				\textbf{\LARGE ATOM}	
			\end{textblock*}
			\begin{textblock*}{0.5\paperwidth}(0.25\paperwidth,0.6\paperheight)
				\centering
				\includegraphics[width=.3\textwidth]{atomic-atom.png}
			\end{textblock*}
		\end{frame}
	}
\end{SectionTitle}


\begin{frame}{The Rutherford–Bohr atomic model}
	\centering
	\only<1>{
	\begin{textblock*}{0.5\paperwidth}(0.25\paperwidth,0.3\paperheight)
		\centering
		\includegraphics[width=\textwidth]{niels-ernst.png}
		Niels Henrik David Bohr  \qquad \qquad Ernest Rutherford
	\end{textblock*}
	}
	\only<2->{	
	\begin{textblock*}{0.45\paperwidth}(0.05\paperwidth,0.2\paperheight)
		\centering
		\includegraphics[width=\textwidth]{bohr-atomic-model-atom.png}
	\end{textblock*}
	\begin{textblock*}{0.45\paperwidth}(0.55\paperwidth,0.2\paperheight)
		\centering
		\begin{itemize}
			\item Protons and Neutrons form the Nucleus
			\item Electrons orbit around 
			\item Atom size is 100000 bigger than the Nucleus size
		\end{itemize}
		\vspace{0.05\textheight}
		\textbf{Is it all?}\\
		\vspace{0.05\textheight}
		\only<3>{
			\textbf{Quantum Mechanics}\\
			\includegraphics[width=0.7\textwidth]{blog_schrodinger1.jpg}
		}
	\end{textblock*}
	}
\end{frame}

\begin{frame}{The quantum mechanical atomic model}
	\includegraphics[width=0.7\textwidth]{Atomic_orbitals.png}
	\begin{textblock*}{0.4\paperwidth}(0.6\paperwidth,0.2\paperheight)
		\centering
		\includegraphics[width=\textwidth]{400px-Hydrogen_transitions.svg.png}
	\end{textblock*}
\end{frame}

\begin{frame}{A closer look}
	\includegraphics[width=0.6\textwidth]{H-d.png}
	\begin{textblock*}{0.4\paperwidth}(0.58\paperwidth,0.25\paperheight)
		\centering
		Different isotopes of the same elements have a slightly different transition wavelenght\\
		\vspace{0.05\textheight}
		\only<2>{
		Result of two different contributions from the nucleus:\\
		\begin{itemize}
			\item Different mass (Mass Shift)
			\item Different size (Field Shift)
		\end{itemize}
		}
	\end{textblock*}
\end{frame}


\begin{frame}{A closer$^2$ look}
	\only<1>{\includegraphics[width=0.6\textwidth]{H-d.png}}
	\only<2>{\hspace{0.1\textwidth}\includegraphics[width=0.4\textwidth]{shapes.png}}
	\begin{textblock*}{0.4\paperwidth}(0.58\paperwidth,0.25\paperheight)
		\centering
		Some isotopes present a further level splitting:\\
		\vspace{0.05\textheight}
		\textbf{HyperFine Structure (HFS)}\\
		\only<2>{
		Result again of two different contributions from the nucleus:\\
		\begin{itemize}
			\item electrons interaction with magnetic moment of nucleus (Magnetic dipole moment)
			\item electrons interaction with electric moment of nucleus (Electric quadrupole moent)
		\end{itemize}
		}
	\end{textblock*}
\end{frame}

\begin{SectionTitle}
	{
		\begin{frame}
			\begin{textblock*}{0.5\paperwidth}(0.25\paperwidth,0.45\paperheight)
				\centering
				\textbf{\LARGE MEASUREMENTS}	
			\end{textblock*}
			\begin{textblock*}{0.5\paperwidth}(0.25\paperwidth,0.6\paperheight)
				\centering
				\includegraphics[width=.3\textwidth]{measuring-tape-tool-svgrepo-com.pdf}
			\end{textblock*}
		\end{frame}
	}
\end{SectionTitle}


\begin{frame}{Orders of magnitude}
	\begin{textblock*}{0.5\paperwidth}(0.25\paperwidth,0.2\paperheight)
		\centering
		\textbf{Optical transition $\rightarrow$} $\sim100$~THz ($10^{14}$~Hz)\\ 
		\vspace{0.05\textheight}
		\textbf{Hyperfine splitting $\rightarrow$} $\sim100$~MHz ($10^{8}$~Hz)\\
	\end{textblock*}
	\only<2>{
	\begin{textblock*}{0.4\paperwidth}(0.05\paperwidth,0.4\paperheight)
		\centering
		\includegraphics[width=0.7\textwidth]{Pub_uni.png}
	\end{textblock*}
	\begin{textblock*}{0.6\paperwidth}(0.4\paperwidth,0.55\paperheight)
		\centering
		Vihre\"{a} Haltijatar $\xleftrightarrow{1.13 km}$ Jyv\"{a}skyl\"{a}n yliopiston fysiikan laitos\\
		We want to know this distance with a better than $\sim$1~mm precision
	\end{textblock*}
	}
\end{frame}

\begin{frame}{Fluorescence}
	\begin{textblock*}{0.45\paperwidth}(0.05\paperwidth,0.3\paperheight)
		\centering
		\includegraphics[width=\textwidth]{fluorescence.pdf}
		\only<1>{\vspace{0.05\textheight}\\Measurement of the number of photons as a function of the laser light wavelenght}
		\only<2>{PROBLEM: Doppler broadening\\ 
		\includegraphics[width=0.4\textwidth]{dopp.png}
		}
	\end{textblock*}
	\begin{textblock*}{0.45\paperwidth}(0.52\paperwidth,0.3\paperheight)
		\centering
		\only<1>{\includegraphics[width=0.9\textwidth]{thermalbroadening-1.pdf}}
		\only<2>{\includegraphics[width=0.9\textwidth]{thermalbroadening-2.pdf}}
	\end{textblock*}
	
\end{frame}

\end{document}
